% SPDX-License-Identifier: MIT
% Copyright (c) 2026 Evgeny Baulin
% Optimization Methods, Seminar 03: one-page cheat sheet (Russian).
% Generated by checks/03. Gradient Descent and Automatic Differentiation.ipynb from web/03/js/content-cheatsheet.js and content/students.json. Do not edit by hand.
% Compile with: latexmk -xelatex -interaction=nonstopmode -halt-on-error Cheatsheet_ru.tex
\RequirePackage{fix-cm}
\documentclass[10pt,a4paper]{article}

\usepackage{amsmath,amssymb}
\usepackage{fontspec}
\setmainfont{cmunrm.otf}[BoldFont=cmunbx.otf, ItalicFont=cmunti.otf, BoldItalicFont=cmunbi.otf]
\setsansfont{cmunss.otf}[BoldFont=cmunsx.otf, ItalicFont=cmunsi.otf, BoldItalicFont=cmunso.otf]
\setmonofont{cmuntt.otf}[BoldFont=cmuntb.otf, ItalicFont=cmunit.otf, BoldItalicFont=cmuntx.otf]
\usepackage{polyglossia}
\setdefaultlanguage{russian}
\setotherlanguage{english}
\newfontfamily\cyrillicfont{cmunrm.otf}[Ligatures=TeX, BoldFont=cmunbx.otf, ItalicFont=cmunti.otf, BoldItalicFont=cmunbi.otf]
\newfontfamily\cyrillicfontsf{cmunss.otf}[Ligatures=TeX, BoldFont=cmunsx.otf, ItalicFont=cmunsi.otf, BoldItalicFont=cmunso.otf]
\newfontfamily\cyrillicfonttt{cmuntt.otf}[BoldFont=cmuntb.otf, ItalicFont=cmunit.otf, BoldItalicFont=cmuntx.otf]
\PolyglossiaSetup{russian}{hyphenmins={2,3}}

\usepackage[a4paper,left=8mm,right=8mm,top=7mm,bottom=11mm,footskip=5.5mm]{geometry}
\usepackage{xcolor}
\usepackage{multicol}
\usepackage{enumitem}
\usepackage{ragged2e}
\usepackage{fancyhdr}
\usepackage[colorlinks=true,linkcolor=hseblue,urlcolor=hseblue]{hyperref}

\definecolor{hseblue}{RGB}{16,45,105}
\definecolor{hseaccent}{RGB}{198,40,40}
\definecolor{hsegray}{RGB}{108,114,126}
\definecolor{hselight}{RGB}{240,243,248}

\hypersetup{pdftitle={Шпаргалка к семинару 3: Градиентный спуск и автоматическое дифференцирование}, pdfauthor={Евгений Баулин}}

% Page style: the author and the license in the footer.
\pagestyle{fancy}
\fancyhf{}
\renewcommand{\headrulewidth}{0pt}
\fancyfoot[L]{\scriptsize\color{hsegray}\textcopyright\ 2026 Евгений Баулин, \href{https://t.me/tarakan_tuc}{t.me/tarakan\_tuc}}
\fancyfoot[R]{\scriptsize\color{hsegray}Лицензия MIT}

% Compact layout.
\setlength{\parindent}{0pt}
\setlength{\parskip}{0pt}
\setlength{\columnsep}{4.5mm}
\setlength{\columnseprule}{0.3pt}
\renewcommand{\columnseprulecolor}{\color{hselight}}
\setlength{\emergencystretch}{2em}
\setlength{\RaggedRightRightskip}{0pt plus 4em}
\binoppenalty=2000
\relpenalty=500
\setlength{\abovedisplayskip}{2pt}
\setlength{\belowdisplayskip}{2pt}
\setlist[itemize]{leftmargin=7pt, labelsep=3pt, itemsep=1.2pt, topsep=1pt, parsep=0pt, partopsep=0pt,
  label={\color{hseaccent}\scriptsize\textbullet}}

\makeatletter
\newcommand{\cssec}[1]{\par\addvspace{3.5pt}\noindent
  \colorbox{hselight}{\parbox{\dimexpr\linewidth-2\fboxsep}{\RaggedRight\hyphenpenalty=10000\color{hseblue}\bfseries\strut #1}}%
  \par\nobreak\vspace{1pt}\@afterheading}
\makeatother
\clubpenalty=10000
\widowpenalty=10000

\begin{document}
\thispagestyle{fancy}

\noindent
\begin{minipage}[b]{0.6\textwidth}
  \RaggedRight\hyphenpenalty=10000
  {\color{hseaccent}\bfseries\small Шпаргалка к семинару 3}\\[1pt]
  {\color{hseblue}\bfseries\Large Градиентный спуск и автоматическое\\ дифференцирование}
\end{minipage}\hfill
\begin{minipage}[b]{0.39\textwidth}
  \raggedleft\footnotesize
  {\bfseries Методы оптимизации}\\
  НИУ ВШЭ, факультет компьютерных наук, 2026--2027\\
  Лектор курса: Андрей Игнатов\\
  Автор: Евгений Баулин, Telegram \href{https://t.me/tarakan_tuc}{@tarakan\_tuc}
\end{minipage}

\vspace{2pt}
{\color{hseblue}\rule{\textwidth}{0.8pt}}
\vspace{-4pt}

\fontsize{9}{10.6}\selectfont
\begin{multicols}{3}
\RaggedRight


% ---- jensen
\cssec{Неравенство Йенсена}
\begin{itemize}
  \item \textbf{Выпуклая} $f$ на выпуклом множестве: $f\big((1-t)x + ty\big) \le (1-t)f(x) + tf(y)$ для всех $x, y$ и $t \in [0, 1]$; хорда лежит не ниже графика, производные не нужны. \textbf{Строго выпуклая}: неравенство строгое при $x \ne y$, $t \in (0, 1)$. Для $x^2$ зазор равен $t(1-t)(x - y)^2$.
  \item \textbf{Три равносильных критерия} (множество открытое и выпуклое): хорда (определение); касательная $f(y) \ge f(x) + \nabla f(x)^\top (y - x)$ для дифференцируемой $f$; кривизна $\nabla^2 f(x) \succeq 0$ для дважды дифференцируемой $f$. Из $\nabla^2 f \succ 0$ всюду следует строгая выпуклость, обратное неверно: $|x|^3$.
  \item \textbf{Веса} $t_i \ge 0$, $\sum_i t_i = 1$: $f\big(\sum_i t_ix_i\big) \le \sum_i t_if(x_i)$; также $f(\mathbb E X) \le \mathbb E f(X)$; для вогнутой функции наоборот: $\log \mathbb E X \ge \mathbb E \log X$. \textbf{Одна пара опровергает} выпуклость: у $x^3$ при $x = -2$, $y = 1$, $t = 0{,}5$ график ($-0{,}125$) выше хорды ($-3{,}5$).
\end{itemize}

% ---- constants
\cssec{Константы $\mu$, $L$ и $\kappa$}
\begin{itemize}
  \item $\mu$-\textbf{сильно выпуклая}: $f - \frac\mu2\|\cdot\|^2$ выпукла, или $\nabla^2 f \succeq \mu I$. $L$-\textbf{гладкая}: $\|\nabla f(x) - \nabla f(y)\| \le L\|x - y\|$, или $-LI \preceq \nabla^2 f \preceq LI$.
  \item \textbf{Две параболы}: $\frac\mu2\|y - x\|^2 \le f(y) - f(x) - \nabla f(x)^\top (y - x) \le \frac L2\|y - x\|^2$. В форме Йенсена зазор между хордой и графиком лежит между $\frac\mu2 t(1-t)\|x - y\|^2$ и $\frac L2 t(1-t)\|x - y\|^2$.
  \item \textbf{Квадратичная} $f = \tfrac{1}{2} x^\top Qx - b^\top x$, $Q \succ 0$: $\nabla f = Qx - b$, $\mu = \lambda_{\min}(Q)$, $L = \lambda_{\max}(Q)$. \textbf{Число обусловленности} $\kappa = L/\mu \ge 1$; большое $\kappa$ — длинная узкая долина.
  \item \textbf{Условие Поляка–Лоясевича} (PL): $\frac{1}{2}\|\nabla f(x)\|^2 \ge \mu\big(f(x) - f^*\big)$ для всех $x$. Оно следует из $\mu$-сильной выпуклости, и при нём любая стационарная точка — точка глобального минимума; выпуклость для него не нужна ($x^2 + 3\sin^2 x$) и не достаточна ($x^4$).
\end{itemize}

% ---- gd
\cssec{Градиентный спуск}
\begin{itemize}
  \item \textbf{Итерация} $x_{k+1} = x_k - \alpha\nabla f(x_k)$ вдоль направления наискорейшего спуска $-\nabla f(x_k)$; $d$ — направление спуска, если $\nabla f(x)^\top d < 0$.
  \item \textbf{Лемма о спуске} ($L$-гладкая $f$): $f(y) \le f(x) + \nabla f(x)^\top (y - x) + \frac L2\|y - x\|^2$. За один шаг $f(x^+) \le f(x) - \alpha\big(1 - \frac{L\alpha}{2}\big)\|\nabla f(x)\|^2$ — убывание при $0 < \alpha < 2/L$; \textbf{безопасный шаг} $1/L$ даёт $f(x^+) \le f(x) - \frac{1}{2L}\|\nabla f(x)\|^2$.
  \item \textbf{Квадратичная} $f$, собственные координаты: $z_{k,i} = (1 - \alpha\lambda_i)^k z_{0,i}$. Сходимость из любой начальной точки $\iff 0 < \alpha < 2/L$. \textbf{Граница устойчивости} $2/L$ не входит: на ней жёсткая координата без конца меняет знак.
  \item \textbf{Коэффициент сжатия} $\rho(\alpha) = \max\big(|1 - \alpha\mu|, |1 - \alpha L|\big)$: $\|x_k - x^*\| \le \rho^k\|x_0 - x^*\|$, $f(x_k) - f^* \le \rho^{2k}\big(f(x_0) - f^*\big)$, $\rho(1/L) = 1 - 1/\kappa$. \textbf{Оптимальный шаг} $\frac{2}{L + \mu}$: $\rho = \frac{\kappa - 1}{\kappa + 1}$.
\end{itemize}

% ---- rates
\cssec{Скорость сходимости}
\begin{itemize}
  \item \textbf{Выпуклая} $f$, $\alpha = 1/L$: $f(x_k) - f^* \le \frac{L\|x_0 - x^*\|^2}{2k} = O(1/k)$, так как сумма неравенств $f(x_{k+1}) - f^* \le \frac L2\big(\|x_k - x^*\|^2 - \|x_{k+1} - x^*\|^2\big)$ телескопируется.
  \item $\mu$-\textbf{сильно выпуклая} $f$ или \textbf{условие PL}, $\alpha = 1/L$: $f(x_k) - f^* \le (1 - \mu/L)^k\big(f(x_0) - f^*\big)$ — из $f(x_{k+1}) \le f(x_k) - \frac{1}{2L}\|\nabla f(x_k)\|^2$ и условия PL, без выпуклости.
  \item \textbf{Число итераций}, чтобы невязка уменьшилась в $1/\varepsilon$ раз: $k \ge \frac{\ln\varepsilon}{\ln(1 - 1/\kappa)}$, примерно $\kappa\ln(1/\varepsilon)$. При $\kappa = 1000$, $\varepsilon = 10^{-6}$: $13809$ при $1/L$ и $3454$ при $\frac{2}{L + \mu}$ на квадратичной функции. Это оценка, а не прогноз: на практике итераций часто меньше.
\end{itemize}

% ---- armijo
\cssec{Бэктрекинг по Армихо}
\begin{itemize}
  \item \textbf{Условие Армихо}: $f(x + \alpha d) \le f(x) + c\alpha\nabla f(x)^\top d$; для $d = -\nabla f$ правая часть равна $f(x) - c\alpha\|\nabla f(x)\|^2$. \textbf{Бэктрекинг}: $\alpha = \alpha_0, \beta\alpha_0, \beta^2\alpha_0, \dots$, пока условие не выполнится; обычно $\beta = \tfrac{1}{2}$, $c = 10^{-4}$.
  \item \textbf{Конечность} ($L$-гладкая $f$): проверку проходит любой $\alpha \le 2(1 - c)/L$, поэтому цикл останавливается и $\alpha \ge \min\big(\alpha_0, 2\beta(1 - c)/L\big)$. Проверяйте $\nabla f(x)^\top d < 0$ и ограничивайте число уменьшений.
\end{itemize}

% ---- stopping
\cssec{Критерии остановки}
\begin{itemize}
  \item \textbf{Масштабированные критерии}: $\|\nabla f(x_k)\| \le \varepsilon_g\max\big(1, |f(x_k)|\big)$, $\|x_{k+1} - x_k\| \le \varepsilon_x\max\big(1, \|x_k\|\big)$, $|f(x_{k+1}) - f(x_k)| \le \varepsilon_f\max\big(1, |f(x_k)|\big)$; в коде $\text{atol} + \text{rtol}\cdot\max(1, s)$.
  \item Критерии по изменению уменьшаются вместе с шагом, $\|x_{k+1} - x_k\| = \alpha_k\|\nabla f(x_k)\|$, и крошечный шаг включает их в любой точке. \textbf{Комбинированное правило}: успех — только по критерию по градиенту; критерий по изменению без него — сигнал застоя; ограничьте число итераций и остановитесь, если $f$ или $\nabla f$ не конечны.
\end{itemize}

% ---- ad
\cssec{Автоматическое дифференцирование}
\begin{itemize}
  \item \textbf{Вычислительный граф}: промежуточные переменные в топологическом порядке; $z = \sin(xy) + x^2$: $u = xy$, $v = \sin u$, $w = x^2$, $z = v + w$. \textbf{Прямой проход} вычисляет и сохраняет значения.
  \item \textbf{Обратный проход}: сопряжённая переменная $\bar v = \partial z/\partial v$; $\bar z = 1$, остальные $0$; вершины обходятся в обратном топологическом порядке, и для каждого ребра $u \to v$ выполняется $\bar u \mathrel{+}= \bar v\,\partial v/\partial u$. Здесь $\bar x = y\cos u + 2x$, $\bar y = x\cos u$.
  \item \textbf{Накопление}: вершина, от которой зависят несколько других, собирает по слагаемому с каждого ребра, $\bar v_i = \sum_{j:\, i \to j} \bar v_j\,\partial v_j/\partial v_i$; для $z = x^2 + x^3$ получаем $\bar x = 2x + 3x^2$. Присваивание сохраняет только последнее слагаемое.
  \item \textbf{Прямой режим} даёт произведение якобиана на вектор $J\dot x$ (JVP): для градиента в $\mathbb R^n$ нужно $n$ проходов. \textbf{Обратный режим} даёт произведение вектора на якобиан $\bar z^\top J$ (VJP): весь градиент скалярной функции за один проход. Проход в небольшое число раз дороже одного вычисления $f$; обратный режим хранит все промежуточные значения (ленту).
\end{itemize}

% ---- fd
\cssec{Конечные разности}
\begin{itemize}
  \item \textbf{Центральная разность} $\frac{f(x + he_i) - f(x - he_i)}{2h}$, ошибка $\approx \frac{h^2}{6}|\partial_i^3 f| + \frac{\varepsilon|f|}{h}$: ошибка отбрасывания плюс ошибка округления, $\varepsilon$ — машинный эпсилон. Лучший $h \sim \varepsilon^{1/3} \approx 6 \cdot 10^{-6}$, ошибка $\sim \varepsilon^{2/3}$. Это проверка градиента, а не способ его вычислять.
\end{itemize}

% ---- logistic
\cssec{Логистическая функция потерь}
\begin{itemize}
  \item \textbf{Производная}: $\frac{d\ell}{ds} = \sigma(s) - y$ для $\ell = \log(1 + e^s) - ys$, $s = w^\top x + b$, $y \in \{0, 1\}$, $\sigma(s) = \frac{1}{1 + e^{-s}}$; поэтому $\nabla_w\ell = \big(\sigma(s) - y\big)x$ и $\frac{\partial\ell}{\partial b} = \sigma(s) - y$.
  \item \textbf{Гладкость} средней функции потерь по строкам $X$: $\nabla^2 f = \frac1n X^\top\operatorname{diag}\big(\sigma_i(1 - \sigma_i)\big)X$ и $\sigma(1 - \sigma) \le \frac{1}{4}$, поэтому $f$ выпукла и $L \le \frac{\lambda_{\max}(X^\top X)}{4n}$; оценка достигается при $w = 0$.
\end{itemize}

% ---- traps
\cssec{Главные ловушки}
\begin{itemize}
  \item При $\alpha = 2/L$ сходимости уже нет: это граница, а не рабочий шаг; предел задаёт $L = \lambda_{\max}$, а $\mu$ влияет только на скорость.
  \item Условие Армихо с $+c\alpha\|\nabla f\|^2$ принимает шаги, увеличивающие $f$; малое изменение $x$ или $f$ может означать крошечный шаг, а не решение.
  \item Обратный проход: $=$ вместо $\mathrel{+}=$ теряет слагаемые; $\sigma(s) - y$ — это $d\ell/ds$, а в градиенте по $w$ есть множитель $x$.
\end{itemize}

\end{multicols}
\end{document}
